\Needspace{7\baselineskip}\Article{303}
The table of negative determinants, arranged by their different classifications, offers many other remarkable observations. For determinants of the form \(-(8n+3)\), the number of classes, both in all the properly primitive genera together and in each separately, is always divisible by \(3\), with the sole exception of the determinant \(-3\). The reason follows directly from Article 256, VI. For determinants whose forms make up a single genus, the number of classes is always odd. For such a determinant there is just one ambiguous class, the principal class; the remaining classes occur in opposite pairs, so their number must be even, and the total odd. This latter property also holds for positive determinants. — Further, the sequence of determinants corresponding to any given classification, that is, to given numbers of genera and classes, always seems to end. We illustrate this rather surprising observation with a few examples. The first number, in Roman numerals, gives the number of properly primitive genera, positive for negative determinants; the next gives the number of classes in each genus. Then follows the sequence of determinants with this classification, their negative sign being omitted for brevity.

\SourceLeaf{372}

\begin{center}\begin{tabularx}{\textwidth}{@{}r@{\quad}X@{}}
I. 1\;\ldots & 1, 2, 3, 4, 7\\[2pt]
I. 3\;\ldots & 11, 19, 23, 27, 31, 43, 67, 163\\[2pt]
I. 5\;\ldots & 47, 79, 103. 127\\[2pt]
I. 7\;\ldots & 71, 151, 223, 343, 463, 487\\[2pt]
II. 1\;\ldots & 5, 6, 8, 9, 10, 12, 13, 15, 16, 18, 22, 25, 28, 37, 58\\[2pt]
II. 2\;\ldots & 14, 17, 20, 32, 34, 36, 39, 46, 49, 52, 55, 63, 64, 73, 82, 97, 100, 142, 148, 193\\[2pt]
IV. 1\;\ldots & 21, 24, 30, 33, 40, 42, 45, 48, 57, 60, 70, 72, 78, 85, 88, 93, 102, 112, 130, 133, 177, 190, 232, 253\\[2pt]
VIII. 1\;\ldots & 105, 120, 165, 168, 210, 240, 273, 280, 312, 330, 345. 357, 385, 408, 462, 520, 760\\[2pt]
XVI. 1\;\ldots & 840, 1320, 1365, 1848\\[2pt]
\end{tabularx}\end{center}
Similarly, there are \(20\) determinants, the largest \(=-1423\), with classification I.~9; \(4\), the largest \(=-1303\), with classification I.~11; and so on. Classifications II.~3; II.~4; II.~5; IV.~2 correspond to no more than \(48,31,44,69\) determinants, respectively, the largest of these being \(-652,-862,-1318\), \(-1012\). The table from which we took these examples has been extended far beyond the largest determinants occurring here\HistoricalNote{\ensuremath{\ast})}{As this is being printed, \textit{continuously} as far as \(-3000\), and also through the entire tenth thousand and several other scattered hundreds, together with very many individual determinants carefully chosen.}, without any further determinants of these classifications appearing. So there seems no doubt that the sequences listed really have ended, and by analogy we may extend the same conclusion to any other classification. For example, since the entire tenth thousand of determinants offered none with fewer than \(24\) classes, it is highly likely that classifications I.~23; I.~21, and so on; II.~11; II.~10, and so on; IV.~5; IV.~4; IV.~3; VIII.~2 had already ended before \(-9000\), or at least apply to very few determinants beyond \(-10000\). But \textit{rigorous} proofs of these observations seem very difficult. — It is no less remarkable that every determinant whose forms divide into \(32\) or more genera has at least two classes in each genus. Thus classifications XXXII.~1, LXIV. 1, and so on, do not occur at all; the smallest such determinant, \(-9240\), has classification XXXII.~2. It seems quite probable that, as the number of genera grows, more and more classifications fail to occur. In this respect the \(65\) determinants given above,
\SourceLeaf{373}
with classifications I.~1; II.~1; IV.~1; VIII.~1; XVI.~1, are very noteworthy. It is easy to see that they, and only they, have both these remarkable properties: all their classes of forms are ambiguous, and any forms contained in the same genus must be both properly and improperly equivalent. These same \(65\) numbers were already given by the distinguished Euler, from a slightly different viewpoint to be mentioned below and with a criterion easy to prove, in \textit{Nouv. Mém. de l’Ac. de Berlin}, 1776, p.\ 338.

\Needspace{7\baselineskip}\Article{304}
The number of properly primitive classes formed by binary forms with a positive \textit{square} determinant \(kk\) can be determined entirely in advance. It equals the number of numbers coprime to \(2k\) and smaller than it. From this, arguments that are not difficult but must be omitted here show that the average number of classes for such determinants around \(kk\) is approximately expressed by \(\frac{8k}{\pi\pi}\). — Positive nonsquare determinants, however, show quite remarkable behavior in this respect. For negative and square determinants, a small number of classes, for example classification I.~1 or I.~3 or II.~1, and so on, occurs only for small determinants and soon ceases altogether. By contrast, by far the greatest part of positive nonsquare determinants, at least those not very large, have classifications with just one class in each genus, so that I.~3; I.~5; II.~2; II.~3; IV.~2, and so on, are very rare. For example, among the \(90\) nonsquare determinants below \(100\), there are \(11,48,27\) with classifications I.~1, II.~1, IV.~1, respectively. Only one, \((37)\), has I.~3; two, \((34\text{ and }82)\), have II.~2; one, \((79)\), has II.~3. Yet as the determinants grow, larger numbers of classes gradually become more frequent. Thus among the \(96\) nonsquare determinants from \(101\) to \(200\), two, \((101,197)\), have I.~3; four, \((145,146,178,194)\), have II.~2; three, \((141,148,189)\), have II.~3. Of the \(197\) determinants from \(801\) to \(1000\), three have I.~3; four II.~2; fourteen II.~3; two II.~5; two II.~6; fifteen IV.~2; six IV.~3; two IV.~4; four VIII.~2; the remaining \(145\) have one class in each genus. — It would be an interesting question, worthy of mathematicians' ingenuity, to investigate the law by which determinants with one class in each genus become steadily rarer. So far we can neither decide theoretically nor conjecture with enough confidence from observation whether they eventually cease altogether, which seems unlikely, or at least become \textit{infinitely rare}, or whether their frequency comes ever closer
\SourceLeaf{374}
to a fixed limit. The average number of classes grows only a little faster than the number of genera, and far more slowly than the square roots of the determinants. Between \(800\) and \(1000\) it is found \(=5,01\). Let us add another observation that partly restores the analogy between positive and negative determinants. We have found that, for a positive determinant \(D\), it is not so much the number of classes itself as this number multiplied by the logarithm of \(t+u\surd D\), where \(t,u\) are the smallest numbers other than \(1,0\) satisfying \(tt-Duu=1\), that is analogous in several ways to the number of classes for a negative determinant. These ways cannot be explained more fully here. The average value of that product is expressed just as accurately by a formula of the form \(m\surd D-n\). But so far we have not been able to determine the values of the constants \(m,n\) theoretically. If any conclusion may be drawn from comparing a few hundreds of determinants, \(m\) seems to differ little from \(2\frac13\). — We reserve for another occasion a fuller discussion of the principles behind the preceding investigations into averages of quantities that do not follow an analytic law, but approach such a law asymptotically ever more closely. We now turn to another investigation, comparing different properly primitive classes of the same determinant, and bringing this long section to an end.

\Needspace{10\baselineskip}\begin{center}\textit{A special algorithm for properly primitive classes: regular and irregular determinants, and related matters.}\end{center}

\Needspace{7\baselineskip}\Article{305}
\textsc{Theorem.} \textit{Let \(K\) be the principal class of forms of a given determinant \(D\), and \(C\) any other class in the principal genus of forms of the same determinant. Let \(2C,3C, 4C\), and so on, be the classes obtained by doubling, tripling, quadrupling, and so on, the class \(C\), as in Article \(249\). If the progression \(C,2C,3C\), and so on, is continued far enough, it eventually reaches a class identical with \(K\). If \(mC\) is the first class identical with \(K\), and the total number of classes in the principal genus \(=n\), then either \(m=n\) or \(m\) is a proper divisor of \(n\).}

\textit{Proof.} I.\ All the classes \(K,C,2C,3C\), and so on, must belong to the principal genus by Article \(247\), so the first \(n+1\) classes of this series, \(K,C,2C\ldots nC\), clearly cannot all be different. Thus either \(K\) is identical with one of the classes \(C,2C,3C\ldots nC\), or at least two of these classes are identical with each other.
\SourceLeaf{375}
Let \(rC=sC\) and \(r>s\). Then also
\[\begin{gathered}(r-1)C=(s-1)C,\quad(r-2)C=(s-2)C\\\text{etc. and }(r+1-s)C=C\end{gathered}\]
giving \((r-s)C=K\). \textit{This proves the first assertion.}

II.\ This also gives at once either \(m=n\) or \(m<n\). It remains only to show that, in the latter case, \(m\) is a proper divisor of \(n\). Since the classes
\[K,C,2C\ldots(m-1)C,\text{ whose collection we denote by }\mathfrak C\]
do not yet exhaust the whole principal genus in this case, let \(C'\) be a class of this genus not contained in \(\mathfrak C\). Consider the collection of classes obtained by composing \(C'\) with each class in \(\mathfrak C\), namely
\[C',C'+C,C'+2C\ldots C'+(m-1)C\text{, denoted by }\mathfrak C'\]
It is easy to see that all classes in \(\mathfrak C'\) differ both from one another and from all those in \(\mathfrak C\), and belong to the principal genus. If \(\mathfrak C\) and \(\mathfrak C'\) exhaust this genus, we therefore have \(n=2m\); if not, \(2m<n\). In the latter case let \(C''\) be a class of the principal genus contained in neither \(\mathfrak C\) nor \(\mathfrak C'\), and consider the collection of classes obtained by composing \(C''\) with each class in \(\mathfrak C\), that is,
\[C'',C''+C,C''+2C\ldots C''+(m-1)C\text{, denoted by }\mathfrak C''\]
These clearly all differ from one another and from all those in \(\mathfrak C\) and \(\mathfrak C'\), and belong to the principal genus. So if \(\mathfrak C,\mathfrak C',\mathfrak C''\) exhaust this genus, \(n=3m\); otherwise \(n>3m\). In that case another class \(C'''\), contained in the principal genus but in neither \(\mathfrak C,\mathfrak C'\) nor \(\mathfrak C''\), can be treated in the same way and shows that either \(n=4m\) or \(n>4m\), and so on. Since \(n\) and \(m\) are finite numbers, the principal genus must eventually be exhausted. Then \(n\) is a multiple of \(m\), or \(m\) a proper divisor of \(n\). \textit{This proves the second assertion.}

\textit{Example.} Let \(D=-356,\ C=(5,2,72)\)\HistoricalNote{\ensuremath{\ast})}{Classes are always expressed here by the simplest \textit{forms} contained in them.}. We find \(2C=(20,8,21), 3C=(4,0,89),\ 4C=(20,-8,21),\ 5C=(5,-2,72),\ 6C=(1,0,356).\) Here
\SourceLeaf{376}
we therefore have \(m=6\), while \(n\) for this determinant is \(12\). Taking for \(C'\) the class \((8,2,45)\), the other five classes in \(\mathfrak C'\) are \((9,-2,40)\), \((9,2,40)\), \((8,-2,45)\), \((17,1,21)\), \((17,-1,21)\).
